% !TEX root = ../main.tex

\documentclass[../main.tex]{subfiles}

\graphicspath{{\subfix{../images/}}}

\begin{document}

\subsection{Калибровка датчиков}

Сырые показания МЭМС-датчиков содержат постоянное смещение нуля (bias),
обусловленное технологическими особенностями изготовления
кристалла~\cite{titterton2004}. Смещения определялись экспериментально по
данным CSV-файла в состоянии покоя. Перевод в физические единицы выполняется
линейным масштабированием:
\begin{equation}
    a_i = \bigl(\tilde{a}_i - b_{a,i}\bigr) \cdot S_a, \quad
    \omega_i = \bigl(\tilde{\omega}_i - b_{\omega,i}\bigr) \cdot S_\omega,
    \label{eq:calib}
\end{equation}
где $\tilde{a}_i$, $\tilde{\omega}_i$~--- сырые целочисленные отсчёты;
$b_{a,i}$, $b_{\omega,i}$~--- смещения нуля;
$S_a = 9{,}81/3050$ [м/с$^2$/отсчёт],
$S_\omega = 0{,}001$ [рад/с/отсчёт].
 
Экспериментально определённые смещения, вынесенные в пространство имён
\texttt{Calib} в файле \texttt{imu\_shared.h}:
\begin{equation}
    \mathbf{b}_a = (31,\; 24,\; 0)^\top, \quad
    \mathbf{b}_\omega = (-7,\; 1{,}5,\; 3)^\top.
    \label{eq:biases}
\end{equation}


\subsection{Представление ориентации кватернионами}

Ориентация тела в пространстве описывается единичным
кватернионом~\cite{kuipers1999}:
\begin{equation}
    \mathbf{q} = q_w + q_x\,\mathbf{i} + q_y\,\mathbf{j} + q_z\,\mathbf{k},
    \quad \|\mathbf{q}\| = 1.
    \label{eq:quat}
\end{equation}
 
Кватернионы предпочтительнее матриц вращения и углов Эйлера: они не страдают
от шарнирного замка (\textit{gimbal lock}), компактны (4~параметра против~9)
и численно стабильны при нормализации~\cite{kuipers1999}.


\subsection{Кватернионное интегрирование гироскопа}
 
Кинематическое уравнение связывает производную кватерниона с угловой скоростью
$\bm{\omega} = (\omega_x,\,\omega_y,\,\omega_z)^\top$ в системе координат
тела~\cite{mahony2008}:
\begin{equation}
    \dot{\mathbf{q}} = \tfrac{1}{2}\,\mathbf{q} \otimes
    \begin{pmatrix} 0 \\ \bm{\omega} \end{pmatrix}.
    \label{eq:qdot}
\end{equation}
 
Для численного интегрирования на шаге $\Delta t$ применяется точная формула
через ось-угол~\cite{titterton2004}:
\begin{equation}
    \theta = \|\bm{\omega}\| \cdot \Delta t, \qquad
    \Delta\mathbf{q} = \Bigl(\cos\tfrac{\theta}{2},\;
    \tfrac{\bm{\omega}}{\|\bm{\omega}\|}\sin\tfrac{\theta}{2}\Bigr).
    \label{eq:dq}
\end{equation}
 
Обновление кватерниона с нормализацией:
\begin{equation}
    \mathbf{q} \leftarrow \frac{\mathbf{q} \otimes \Delta\mathbf{q}}
    {\|\mathbf{q} \otimes \Delta\mathbf{q}\|}.
    \label{eq:qupdate}
\end{equation}
 
\subsection{Матрица поворота из кватерниона}
 
Для перевода вектора ускорения из системы координат тела в мировую используется
матрица вращения $\mathbf{R} \in SO(3)$~\cite{kuipers1999}:
\begin{equation}
    \mathbf{R} =
    \begin{pmatrix}
        1{-}2(q_y^2{+}q_z^2) & 2(q_xq_y{-}q_wq_z) & 2(q_xq_z{+}q_wq_y) \\
        2(q_xq_y{+}q_wq_z) & 1{-}2(q_x^2{+}q_z^2) & 2(q_yq_z{-}q_wq_x) \\
        2(q_xq_z{-}q_wq_y) & 2(q_yq_z{+}q_wq_x) & 1{-}2(q_x^2{+}q_y^2)
    \end{pmatrix}.
    \label{eq:rot}
\end{equation}
 
\subsection{Углы Эйлера (конвенция ZYX)}
 
Из кватерниона вычисляются углы крена $\phi$, тангажа $\vartheta$ и рыскания
$\psi$~\cite{titterton2004}:
\begin{align}
    \phi      &= \operatorname{atan2}\!\bigl(2(q_wq_x+q_yq_z),\;
                  1-2(q_x^2+q_y^2)\bigr), \label{eq:roll}\\
    \vartheta &= \arcsin\!\bigl(2(q_wq_y-q_zq_x)\bigr), \label{eq:pitch}\\
    \psi      &= \operatorname{atan2}\!\bigl(2(q_wq_z+q_xq_y),\;
                  1-2(q_y^2+q_z^2)\bigr). \label{eq:yaw}
\end{align}
 
\subsection{Двойное интегрирование (скорость и позиция)}
 
Ускорение переводится из системы тела в мировую, из него вычитается
гравитация~\cite{groves2013}:
\begin{equation}
    \mathbf{a}_\mathrm{world} = \mathbf{R}\,\mathbf{a}_\mathrm{body}
    - \mathbf{g}, \quad
    \mathbf{g} = (0,\;0,\;{-}9{,}81)^\top \;\text{[м/с}^2\text{]}.
    \label{eq:aworld}
\end{equation}
 
Интегрирование методом Эйлера первого порядка:
\begin{align}
    \mathbf{v}(t{+}\Delta t) &= \mathbf{v}(t) +
        \mathbf{a}_\mathrm{world}(t)\cdot\Delta t, \label{eq:vel}\\
    \mathbf{p}(t{+}\Delta t) &= \mathbf{p}(t) +
        \mathbf{v}(t)\cdot\Delta t. \label{eq:pos}
\end{align}
 
\subsection{Метод ZUPT}
 
ZUPT (\textit{Zero Velocity Update})~--- классический способ частичной
коррекции в моменты покоя~\cite{skog2010}. Если норма угловой скорости
опускается ниже порогового значения $\omega_\mathrm{th}=0{,}015$~рад/с,
скорость мягко гасится с коэффициентом $\alpha=0{,}85$:
\begin{equation}
    \text{если}\ \|\bm{\omega}\| < \omega_\mathrm{th}:\quad
    \mathbf{v} \leftarrow \mathbf{v}\cdot\alpha.
    \label{eq:zupt}
\end{equation}
Дополнительно, если компонента мирового ускорения по модулю меньше
$0{,}3$~м/с$^2$, она обнуляется полностью, что устраняет интегрирование
шумового сигнала акселерометра в статике.

	
\end{document}
